\documentclass[12pt,a4paper]{article}
\usepackage[utf8]{inputenc}
\usepackage[T1]{fontenc}
\usepackage{amsmath,amssymb,amsthm}
\usepackage{physics}
\usepackage{geometry}
\usepackage{enumitem}
\usepackage{hyperref}
\usepackage{fancyhdr}
\usepackage{titlesec}

\geometry{margin=1in}
\pagestyle{fancy}
\fancyhf{}
\rhead{Lecture 9}
\lhead{Advanced Quantum Mechanics}
\cfoot{\thepage}

\titleformat{\section}{\large\bfseries}{}{0em}{}
\titleformat{\subsection}{\normalsize\bfseries}{}{0em}{}

\newtheorem{definition}{Definition}

\title{\textbf{Lecture 9: Interactions and the Dirac Equation} \\ \large Momentum Conservation, Feynman Diagrams, and Relativistic Electrons}
\author{Based on Leonard Susskind's \textit{Advanced Quantum Mechanics}}
\date{}

\begin{document}
\maketitle
\thispagestyle{fancy}

% ============================================================
\section{1. Momentum Conservation in Quantum Field Theory}
% ============================================================

\subsection{1.1 The Hamiltonian in momentum space}

Consider the second-quantized Hamiltonian for non-interacting particles with a constant potential $mc^2$:
\[
    H = \int dx\;\Psi^\dagger(x)\left[-\frac{\hbar^2}{2m}\frac{d^2}{dx^2} + mc^2\right]\Psi(x).
\]
Expanding the field operators in the momentum basis via Fourier transform:
\[
    \Psi(x) = \int\frac{dp}{\sqrt{2\pi\hbar}}\;\tilde{a}(p)\,e^{ipx/\hbar}, \qquad
    \Psi^\dagger(x) = \int\frac{dq}{\sqrt{2\pi\hbar}}\;\tilde{a}^\dagger(q)\,e^{-iqx/\hbar}.
\]

\subsection{1.2 The delta function and momentum conservation}

Substituting and integrating over $x$:
\[
    \int dx\;e^{i(p - q)x/\hbar} = 2\pi\hbar\;\delta(p - q).
\]
The delta function enforces $p = q$: the momentum of the annihilated particle must equal the momentum of the created particle.  The result:
\[
    H_{\text{mass}} = mc^2\int dp\;\tilde{a}^\dagger(p)\,\tilde{a}(p),
\]
which simply counts each particle and assigns it an energy $mc^2$.  The kinetic term similarly gives:
\[
    H_{\text{kin}} = \int dp\;\frac{p^2}{2m}\;\tilde{a}^\dagger(p)\,\tilde{a}(p).
\]

\subsection{1.3 General momentum conservation rule}

For any Hamiltonian of the form $H = \int dx\;\mathcal{H}[\Psi^\dagger(x), \Psi(x)]$---where the integrand involves products of field operators all evaluated at the \emph{same} point $x$ and integrated over all space---the integration over $x$ produces a delta function:
\[
    \delta\!\left(\sum_{\text{created}} q_i - \sum_{\text{annihilated}} p_j\right).
\]
This enforces \textbf{total momentum conservation}: the sum of momenta of particles created must equal the sum of momenta of particles annihilated.  Individual momenta may change, but the total is preserved.

% ============================================================
\section{2. Interaction Terms}
% ============================================================

\subsection{2.1 Contact scattering}

Consider two species of particles (e.g., ``electrons'' and ``protons'').  A short-range interaction at the same point:
\[
    H_{\text{int}} = g\int dx\;\Psi_e^\dagger(x)\,\Psi_p^\dagger(x)\,\Psi_e(x)\,\Psi_p(x).
\]
This operator:
\begin{enumerate}[nosep]
    \item Finds an electron and a proton at the same position $x$.
    \item Annihilates both.
    \item Creates an electron and a proton at the same point (possibly with different momenta).
\end{enumerate}
The coupling constant $g$ determines the strength of the interaction.  Total momentum is conserved, but individual momenta can change---this is \textbf{scattering}.

\subsection{2.2 Particle decay}

If particle $A$ can decay into particles $B$ and $C$:
\[
    H_{\text{decay}} = g\int dx\;\bigl[\Psi_B^\dagger(x)\,\Psi_C^\dagger(x)\,\Psi_A(x) + \Psi_A^\dagger(x)\,\Psi_C(x)\,\Psi_B(x)\bigr].
\]
The first term annihilates an $A$ and creates a $B$ and $C$.  The second term (the Hermitian conjugate) describes the reverse process: $B + C \to A$.  The Hamiltonian must be Hermitian, so both directions are always present.

\subsection{2.3 Coupling constants}

The parameter $g$ is the \textbf{coupling constant}:
\begin{itemize}[nosep]
    \item Small $g$: weak interaction, low scattering probability.
    \item Large $g$: strong interaction, high scattering probability.
    \item $g$ is typically determined experimentally.
\end{itemize}

% ============================================================
\section{3. Feynman Diagrams}
% ============================================================

\subsection{3.1 Iterating the Hamiltonian}

The time-evolution operator expanded to second order:
\[
    U(t) \approx 1 - i\epsilon H - \frac{\epsilon^2}{2}H^2 + \cdots
\]
The first-order term $H$ gives direct processes (single vertex).  The second-order term $H^2$ gives compound processes (two vertices), leading to a rich set of scattering and exchange phenomena.

\subsection{3.2 Exchange diagrams}

If the Hamiltonian contains a term $A \to B + C$, then $H^2$ contains the process:
\[
    A + B \;\longrightarrow\; C \;\longrightarrow\; A + B,
\]
where $C$ is a temporary (virtual) intermediate particle.  This is an \textbf{exchange diagram}: the interaction between $A$ and $B$ is mediated by the exchange of $C$.  The amplitude carries a factor $g^2$.

\subsection{3.3 Self-energy corrections}

Another second-order process: a particle $A$ emits a $B$ and $C$, which then recombine into $A$:
\[
    A \;\longrightarrow\; B + C \;\longrightarrow\; A.
\]
This modifies the effective properties (mass, charge) of particle $A$.  Example: an electron emitting and reabsorbing a virtual photon---the electron has approximately a 1\% probability of containing a virtual photon at any given time.

\subsection{3.4 Vertices and diagrams}

Each vertex in a Feynman diagram corresponds to a term in the Hamiltonian and carries one factor of the coupling constant $g$.  A diagram with $n$ vertices contributes at order $g^n$.  The full amplitude for any process is the sum over all possible diagrams, providing a systematic perturbative expansion.

% ============================================================
\section{4. Introduction to the Dirac Equation}
% ============================================================

\subsection{4.1 Motivation}

For a relativistic particle, the energy--momentum relation is:
\[
    E^2 = p^2c^2 + m^2c^4.
\]
Dirac wanted a wave equation that is \emph{first order} in time (like the Schr\"odinger equation) rather than second order (like the Klein--Gordon equation $-\hbar^2\partial_t^2\psi = (-\hbar^2c^2\partial_x^2 + m^2c^4)\psi$).

\subsection{4.2 The one-dimensional massless case}

For a massless particle moving to the right, $E = p$ (setting $c = 1$):
\[
    i\hbar\frac{\partial\psi}{\partial t} = -i\hbar\frac{\partial\psi}{\partial x}.
\]
Solutions are right-moving waves $\psi(x - t)$.  Similarly, $E = -p$ gives left-moving waves.

\subsection{4.3 Two-component structure}

To describe both left- and right-movers, introduce a two-component object $\Psi = \begin{pmatrix} \psi_1 \\ \psi_2 \end{pmatrix}$ and a matrix $\alpha$:
\[
    H = \alpha\, p, \qquad \alpha = \begin{pmatrix} 1 & 0 \\ 0 & -1 \end{pmatrix} = \sigma_z.
\]
The eigenvalue $+1$ corresponds to right-movers and $-1$ to left-movers.

\subsection{4.4 Adding mass: the $\beta$ matrix}

To obtain $H^2 = p^2 + m^2$, add a mass term:
\[
    H = \alpha\, p + \beta\, m.
\]
Squaring: $H^2 = \alpha^2 p^2 + \beta^2 m^2 + (\alpha\beta + \beta\alpha)pm$.  We need:
\begin{align}
    \alpha^2 &= 1, \qquad \beta^2 = 1, \qquad \alpha\beta + \beta\alpha = 0.
\end{align}
The solution uses two anti-commuting Pauli matrices:
\[
    \alpha = \sigma_z = \begin{pmatrix} 1 & 0 \\ 0 & -1 \end{pmatrix}, \qquad
    \beta = \sigma_x = \begin{pmatrix} 0 & 1 \\ 1 & 0 \end{pmatrix}.
\]

\subsection{4.5 The one-dimensional Dirac equation}

The full equation:
\[
    i\hbar\frac{\partial}{\partial t}\begin{pmatrix}\psi_1 \\ \psi_2\end{pmatrix} = \left[\alpha\!\left(-i\hbar\frac{\partial}{\partial x}\right) + \beta\, m\right]\begin{pmatrix}\psi_1 \\ \psi_2\end{pmatrix}.
\]
Written out:
\begin{align}
    i\hbar\,\dot{\psi}_1 &= -i\hbar\,\psi_1' + m\,\psi_2, \\
    i\hbar\,\dot{\psi}_2 &= +i\hbar\,\psi_2' + m\,\psi_1.
\end{align}
The mass term $\beta m$ \emph{couples} the two components: it mixes left-movers and right-movers, giving the particle a mass and allowing it to move slower than the speed of light.

\medskip
\noindent\textit{The one-dimensional Dirac equation captures the essential physics: the necessity of a multi-component wave function, the role of the mass as a coupling between chiralities, and the emergence of both positive and negative energy solutions.  In the next lecture, we extend this to three dimensions and confront the Dirac sea.}

\end{document}
